\phantomsection  % so hyperref creates bookmarks
\addcontentsline{toc}{chapter}{Información sobre Licencia}
%\label{label_fdl}

 \begin{center}

       Información sobre Licencia


\end{center}

\section*{Licencia de la documentación}

\todo{Confirmar con el tutor la licencia de la memoria (propuesta: CC BY-SA 4.0).}
Esta memoria se distribuye bajo la licencia \textit{Creative Commons
Reconocimiento-CompartirIgual 4.0 Internacional} (CC BY-SA
4.0)~\cite{cc-by-sa-4}. Se permite copiar, redistribuir y adaptar el
documento con cualquier finalidad, siempre que se reconozca la autoría, se
indiquen los cambios realizados y las obras derivadas se distribuyan con la
misma licencia.

Esta licencia se aplica únicamente al texto, las tablas y las figuras
originales de la memoria. No se extiende al software descrito en ella, a los
fragmentos de código que se reproducen a título ilustrativo ni a las marcas y
logotipos de terceros que puedan aparecer, que conservan sus propias
condiciones.

\section*{Licencia del software}

% PENDIENTE (P-06): forma de presentar las bases de partida, a confirmar con el tutor.
PymeKit no es un desarrollo íntegramente original. Como se explica en la
introducción y en el capítulo de antecedentes, el software parte de dos bases
de código de terceros distribuidas con \textbf{licencia comercial} ---una base
de aplicación SaaS y un CMS para la base de datos---, que se han integrado,
adaptado y extendido en este trabajo. El autor adquirió ambas con su licencia
correspondiente, que permite usarlas como base de productos propios pero no
publicar ni redistribuir su código fuente.
\todo{Confirmar este resumen con el texto de las licencias adquiridas.}

Por ese motivo:

\begin{itemize}
  \item el repositorio del proyecto es \textbf{privado} y su código fuente
        \textbf{no se publica ni se redistribuye} bajo ninguna licencia
        libre;
  \item el código se pone a disposición del tribunal y del tutor únicamente
        para la evaluación del Trabajo Fin de Grado;
        \todo{Confirmar con el tutor el mecanismo de entrega del código al
        tribunal (acceso al repositorio privado o copia en soporte digital).}
  \item el origen del código reutilizado está documentado internamente en el
        repositorio, junto con el registro de los módulos portados y de las
        modificaciones realizadas sobre ellos; y
  \item las aportaciones propias del trabajo se integran en ese mismo
        repositorio y quedan sujetas a sus mismas condiciones de distribución.
\end{itemize}

Las dependencias de código abierto que usa el proyecto (por ejemplo, React,
TanStack, Supabase, Hono, Drizzle, Tailwind CSS o Playwright) se rigen por sus
respectivas licencias, recogidas en sus paquetes.
